% !TeX program = pdflatex
% !TeX encoding = UTF-8
% Preamble.
% !TeX program = pdflatex
% !TeX encoding = UTF-8
% Preamble.
\documentclass[12pt]{article}

% Packages.
\usepackage[utf8]{inputenc}     % Unicode support (Umlauts etc.)
\usepackage{amsmath, amssymb}   % Advanced math typesetting
\usepackage{stmaryrd, bbm}
\usepackage[french]{babel}      % Change hyphenation rules
\usepackage[T1]{fontenc}
\usepackage{geometry}
\usepackage{listings}           % Source code formatting and highlighting
\usepackage{csquotes}
\usepackage{interval}
\usepackage{pifont}
\usepackage{import}

\import{../../}{Common_MP2I}
\date{today{}}

\begin{document}

\paragraph{Question de cours.}
Soit $f : \R \to \R$ qui réalise une bijection.
Lien entre la monotonie de $f$ et de $f^{-1}$. Le démontrer.

\paragraph{Question 1.} 
Soit $f : x \mapsto \sqrt{x^2 + x + 6}$
Admet-elle un asymptote en $+\infty$ ? Si oui, la déterminer.

\paragraph{Question 2.}
Montrer que $\forall x \in \R_{+}, \forall n \in \mathbb{N}$, $$
    e^x \ge \sum_{k=0}^{n} \frac{x^k}{k!}
$$

Que se passe-t-il pour $x \in \R$ ?

\paragraph{Question 3.} 
Soit $f : x \mapsto \frac{x}{1 + \abs{x}}$
Montrer que $f$ réalise une bijection de $\R$ sur un ensemble
$J$ à préciser.

\end{document}
